% Figure 8: How to orient a boundary
\begin{tikzpicture}[line cap=round,line join=round,>=Stealth,scale=0.9,
    midarrow/.style={postaction={decorate,
      decoration={markings,mark=at position #1 with {\arrow{Stealth}}}}}]
  % left: an oriented 1-manifold with boundary points -1 and +1
  \draw[thick,midarrow=0.25,midarrow=0.75]
    (-3.4,-1.1) .. controls (-3.0,-0.5) and (-3.5,0.0) .. (-3.15,0.45)
    .. controls (-2.9,0.75) and (-2.55,0.6) .. (-2.6,0.3)
    .. controls (-2.65,0.05) and (-3.0,0.15) .. (-2.9,0.5)
    .. controls (-2.8,0.85) and (-2.5,1.05) .. (-2.15,1.2);
  \fill (-3.4,-1.1) circle (1.7pt);
  \node[below] at (-3.4,-1.18) {$-1$};
  \fill (-2.15,1.2) circle (1.7pt);
  \node[above right] at (-2.15,1.22) {$+1$};
  % right: a 2-manifold with boundary, with basis (v1,v2) at a boundary point
  \fill[pattern=dots,pattern color=gray,even odd rule]
    plot[smooth cycle,tension=0.9] coordinates
      {(-0.4,-0.4) (0.6,-1.15) (1.9,-1.0) (2.6,-0.1) (2.1,0.95) (0.8,1.25) (-0.15,0.6)}
    (1.15,0.25) circle (0.42);
  \draw[thick] plot[smooth cycle,tension=0.9] coordinates
    {(-0.4,-0.4) (0.6,-1.15) (1.9,-1.0) (2.6,-0.1) (2.1,0.95) (0.8,1.25) (-0.15,0.6)};
  \draw[thick] (1.15,0.25) circle (0.42);
  % outward vector v1 and tangent vector v2 at a boundary point of the hole
  \fill (1.15,-0.17) circle (1.4pt);
  \draw[->] (1.15,-0.17) -- (1.15,-0.62);
  \node at (1.38,-0.62) {\footnotesize $v_1$};
  \draw[->] (1.15,-0.17) -- (0.68,-0.12);
  \node at (0.62,0.08) {\footnotesize $v_2$};
  % small tangent arrow on the outer boundary
  \draw[->] (2.35,0.62) -- (2.6,0.32);
\end{tikzpicture}
